% Part: incompleteness
% Chapter: representability-in-q
% Section: sigma1-completeness

\documentclass[../../../include/open-logic-section]{subfiles}

\begin{document}

\olfileid[pt]{inc}{inp}{s1c}
\olsection{Completude \texorpdfstring{$\Sigma_1$}{Sigma-1}}

Apesar da incompletude de $\Th{Q}$ e de suas extensões consistentes e
axiomatizáveis, vimos que $\Th{Q}$ de fato prova muitos fatos básicos
sobre numerais. Na verdade, isso pode ser ampliado consideravelmente.
Para compreender o alcance do que pode ser provado em~$\Th{Q}$,
introduzimos as noções de !!{formula}s $\Delta_0$, $\Sigma_1$ e
$\Pi_1$. Em termos aproximados, !!a{formula} $\Sigma_1$ tem a forma
$\lexists[x][!B(x)]$, em que $!B$ é construída usando apenas conectivos
proposicionais e quantificadores limitados. Mostraremos que, se $!A$ é
!!a{sentence} $\Sigma_1$ verdadeira em $\Struct{N}$, então
$\Th{Q} \Proves !A$
(\olref{thm:sigma1-completeness}).

\begin{defn}
\ollabel{defn:bd-quant}
Uma \emph{!!{formula} existencial limitada} tem a forma
$\lexists[x][(x < t \land !A(x))]$, em que $t$ é um termo qualquer, e
por convenção a escrevemos como $\bexists{x < t}{!A(x)}$.
%
Uma \emph{!!{formula} universal limitada} tem a forma
$\lforall[x][(x < t \lif !A(x))]$, em que $t$ é um termo qualquer, e
por convenção a escrevemos como $\bforall{x < t}{!A(x)}$.
\end{defn}

\begin{defn}
\ollabel{defn:delta0-sigma1-pi1-frm}
Uma !!{formula} $!B$ é $\Delta_0$ se for construída a partir de
!!{formula}s atômicas usando apenas conectivos proposicionais e
quantificação limitada.
%
Uma !!{formula} $!A$ é $\Sigma_1$ se
$!A \ident \lexists[x][!B(x)]$, em que $!B$ é $\Delta_0$.
%
Uma !!{formula} $!A$ é $\Pi_1$ se $!A \ident \lforall[x][!B(x)]$, em
que $!B$ é $\Delta_0$.
\end{defn}

\begin{lem}
\ollabel{lem:q-proves-clterm-id} Suponha que $t$ seja um termo fechado
tal que $\Value{t}{N} = n$. Então $\Th{Q} \Proves \eq[t][\num n]$.
\end{lem}

\begin{proof}
Provamos isso por indução sobre a complexidade de~$t$. No caso base,
$\Value{\Obj 0}{N} = 0$, e $\Th{Q} \Proves \eq[\Obj 0][\num 0]$, pois
$\num 0 \ident \Obj 0$.
%
Para o caso indutivo, sejam $t_1$ e $t_2$ termos tais que
$\Value{t_1}{N} = n_1$, $\Value{t_2}{N} = n_2$,
$\Th{Q} \Proves \eq[t_1][\num n_1]$, e
$\Th{Q} \Proves \eq[t_2][\num n_2]$.

Então $\Value{(t_1')}{N} = n_1 + 1$, e temos que $\Th{Q} \Proves
\eq[t_1'][{\num n_1}']$ pelas regras de primeira ordem para a identidade,
aplicadas à hipótese de indução e à !!{formula}
$\eq[\num{n_1}'][\num{n_1}']$,
de modo que temos $\Th{Q} \Proves \eq[t_1'][\num{n_1 + 1}]$ pela
definição dos numerais.

Para somas, temos
\[
      \Value{(t_1 + t_2)}{N}
    = \Value{t_1}{N} + \Value{t_2}{N}
    = n_1 + n_2.
\]
Pela hipótese de indução e pelas regras da identidade,
$\Th{Q} \Proves \eq[t_1 + t_2][\num{n_1} + t_2]$, e então
$\Th{Q} \Proves \eq[t_1 + t_2][\num{n_1} + \num{n_2}]$
por uma segunda aplicação das regras da identidade. Por
\olref[inc][req][bre]{lem:q-proves-add},
$\Th{Q} \Proves \eq[\num{n_1} + \num{n_2}][\num{n_1 + n_2}]$,
portanto $\Th{Q} \Proves \eq[t_1 + t_2][\num{n_1 + n_2}]$.

Um raciocínio semelhante também funciona para~$\times$, usando
\olref[inc][req][bre]{lem:q-proves-mult}.
%
Como isso esgota os termos fechados da aritmética, temos que
$\Th{Q} \Proves \eq[t][\num n]$ para todo termo fechado~$t$ tal que
$\Value{t}{N} = n$.
\end{proof}

\begin{prob}
Prove em detalhe a parte de \olref{lem:q-proves-clterm-id} que
envolve~$\times$.
\end{prob}

\begin{lem}
\ollabel{lem:atomic-completeness}
Suponha que $t_1$ e $t_2$ sejam termos fechados. Então
\begin{enumerate}
\item Se $\Value{t_1}{N} = \Value{t_2}{N}$, então
    $\Th{Q} \Proves \eq[t_1][t_2]$.
\item Se $\Value{t_1}{N} \neq \Value{t_2}{N}$, então
    $\Th{Q} \Proves \eq/[t_1][t_2]$.
\item Se $\Value{t_1}{N} < \Value{t_2}{N}$, então
    $\Th{Q} \Proves t_1 < t_2$.
\item Se $\Value{t_2}{N} \leq \Value{t_1}{N}$, então
    $\Th{Q} \Proves \lnot(t_1 < t_2)$.
\end{enumerate}
\end{lem}

\begin{proof}
Dados os termos $t_1$ e $t_2$, fixamos $n = \Value{t_1}{N}$ e
$m = \Value{t_2}{N}$.

Suponha $!A \ident t_1 = t_2$. Por \olref{lem:q-proves-clterm-id},
$\Th{Q} \Proves \eq[t_1][\num n]$ e $\Th{Q} \Proves \eq[t_2][\num n]$.
Se $n = m$, então $\Th{Q} \Proves \eq[\num n][\num m]$ e, portanto,
$\Th{Q} \Proves \eq[t_1][t_2]$ pela transitividade da identidade.
Se $n \neq m$, então $\Th{Q} \Proves \eq/[\num n][\num m]$ e,
novamente pela transitividade da identidade,
$\Th{Q} \Proves \eq/[t_1][t_2]$.

Seja agora $!A \ident t_1 < t_2$. Em ambos os casos, usamos o
axioma~$!Q_8$, que afirma $x < y \liff \lexists[z][\eq[z' + x][y]]$
para todos $x,y$.

Suponha $\Sat{N}{t_1 < t_2}$. Então existe algum $k \in \Nat$ tal que
$n + k + 1 = m$. Por \olref{lem:q-proves-clterm-id},
$\Th{Q} \Proves \eq[t_1][\num n]$ e $\Th{Q} \Proves \eq[t_2][\num m]$
e, pela primeira parte deste lema,
$\Th{Q} \Proves \eq[\num n + {\num k}'][\num m]$.
Pela transitividade da identidade, segue que
$\Th{Q} \Proves \eq[{\num k}' + t_1][t_2]$,
portanto $\Th{Q} \Proves \lexists[z][\eq[z' + t_1][t_2]]$.
Pela direção da direita para a esquerda de~$!Q_8$,
$\Th{Q} \Proves t_1 < t_2$.

Suponha, em vez disso, que $\Sat/{N}{t_1 < t_2}$, isto é, $m \leq n$.
%
Trabalhamos em~$\Th{Q}$ e supomos que $t_1 < t_2$. Pela direção da
esquerda para a direita de~$!Q_8$, existe algum~$z$ tal que
$\eq[z' + t_1][t_2]$. Como $\Th{Q} \Proves \eq[t_1][\num n]$ e
$\Th{Q} \Proves \eq[t_2][\num m]$, $\eq[z' + \num n][\num m]$.
%
Por uma indução externa sobre~$m$, usando~$!Q_5$,
$\eq[z' + \num{n - m}][\Obj 0]$.
Se $m = n$, então $\eq/[z'][\Obj 0]$, o que produz uma contradição
por~$!Q_3$. Se $m < n$, então
$\eq[(z' + \num{n - m - 1})'][\Obj 0]$ novamente por~$!Q_5$, o que
produz uma contradição por~$!Q_3$. Logo,
$\Th{Q} \Proves \lnot(t_1 < t_2)$.
\end{proof}

\begin{lem}
\ollabel{lem:bounded-quant-equiv}
Suponha que $!A$ seja !!a{formula}, $t$ um termo fechado e
$k=\Value{t}{N}$. Então
\begin{enumerate}
\item $\Th{Q} \Proves \bforall{x<t}{!A(x)}$ sse $\Th{Q} \Proves
    !A(\num 0) \land \dots \land !A(\num{k-1})$.
\item $\Th{Q} \Proves \bexists{x<t}{!A(x)}$ sse $\Th{Q} \Proves
    !A(\num 0) \lor \dots \lor !A(\num{k-1})$.
\end{enumerate}
\end{lem}

\begin{proof}
    Provamos o caso do quantificador universal limitado. Se
    $\Value{t}{N} = 0$, então o lado esquerdo da equivalência é
    demonstrável em~$\Th{Q}$, pois não existe $x<\num 0$ por
    \olref[inc][req][min]{lem:less-zero}. De modo semelhante, podemos
    tomar uma conjunção vazia simplesmente como $\ltrue$, que também é
    demonstrável em~$\Th{Q}$.
    %
    Supomos, portanto, que $\Value{t}{N} = k+1$ para algum número
    natural~$k$. Por \olref{lem:q-proves-clterm-id}, podemos supor que
    estamos trabalhando com !!a{formula} da forma
    $\bforall{x<\num{k+1}}{!A(x)}$.
    
    Suponha que $\Th{Q} \Proves \bforall{x<\num{k+1}}{!A(x)}$ e seja
    $n \leq k$. Como $\Th{Q} \Proves \num n < \num{k+1}$ por
    \olref{lem:atomic-completeness}, segue pela lógica que
    $\Th{Q} \Proves !A(\num n)$. Aplicando esse fato $k+1$ vezes, uma
    para cada $n \leq k$, obtemos
    $\Th{Q} \Proves !A(\num 0) \land \dots \land !A(\num k)$, como
    desejado.
    
    Para a outra direção, suponha que
    $\Th{Q} \Proves !A(\num 0) \land \dots \land !A(\num k)$.
    Trabalhando em $\Th{Q}$, suponha que $x < \num{k+1}$. Por
    \olref[inc][req][min]{lem:less-nsucc}, temos
    $x = \num 0 \lor \dots \lor x = \num k$; pela lógica, segue
    então~$!A(x)$ e, portanto, segue a afirmação universal
    $\bforall{x<\num{k+1}}{!A(x)}$.
    
    A prova da equivalência para !!{formula}s quantificadas
    existencialmente de modo limitado é semelhante.
\end{proof}

\begin{prob}
Apresente uma prova detalhada do caso existencial em
\olref{lem:bounded-quant-equiv}.
\end{prob}

\begin{lem}
\ollabel{lem:delta0-completeness}
Se $!A$ é !!a{sentence} $\Delta_0$ verdadeira em $\Struct{N}$, então
$\Th{Q} \Proves !A$.
\end{lem}

\begin{proof}
Provamos isso por indução sobre a complexidade da !!{formula}.
%
O caso base é dado por \olref{lem:atomic-completeness}; passamos,
portanto, ao passo de indução. Para simplificar, dividimos o caso da
negação em subcasos, dependendo da estrutura da !!{formula} à qual ela
é aplicada.

\begin{enumerate}
\item Suponha que $(!A \land !B)$ seja verdadeira em $\Struct{N}$;
então $!A$ e $!B$ são verdadeiras em~$\Struct{N}$. Pela hipótese de
indução, $\Th{Q} \Proves !A$ e $\Th{Q} \Proves !B$; logo,
$\Th{Q} \Proves (!A \land !B)$ pela lógica.
%
\item Suponha que $\lnot (!A \land !B)$ seja verdadeira em
$\Struct{N}$; então ou $\lnot !A$ ou $\lnot !B$ é verdadeira em
$\Struct{N}$. Sem perda de generalidade, suponha a primeira. Pela
hipótese de indução, $\Th{Q} \Proves \lnot !A$ e, portanto,
$\Th{Q} \Proves \lnot (!A \land !B)$ pela lógica.
%
\item Suponha que $(!A \lor !B)$ seja verdadeira em $\Struct{N}$;
então ou $!A$ é verdadeira em $\Struct{N}$ ou $!B$ é verdadeira em
$\Struct{N}$. Sem perda de generalidade, suponha a primeira. Pela
hipótese de indução, $\Th{Q} \Proves !A$ e, portanto,
$\Th{Q} \Proves (!A \lor !B)$ pela lógica.
%
\item Suponha que $\lnot(!A \lor !B)$ seja verdadeira em
$\Struct{N}$; então $\lnot !A$ e $\lnot !B$ são verdadeiras em
$\Struct{N}$. Pela hipótese de indução,
$\Th{Q} \Proves \lnot !A$ e $\Th{Q} \Proves \lnot !B$. Por
conseguinte, $\Th{Q} \Proves \lnot(!A \lor !B)$ pela lógica.
%
\item Suponha que $\bforall{x<t}{!A(x)}$ seja verdadeira
em~$\Struct{N}$, em que $t$ é um termo fechado e $k=\Value{t}{N}$.
Pela hipótese de indução e pela lógica, se $!A(\num n)$ é verdadeira
em~$\Struct{N}$ para todo $n < \Value{t}{N}$, então
$\Th{Q} \Proves !A(\num 0) \land \dots \land !A(\num{k-1})$. Por
\olref{lem:bounded-quant-equiv}, segue que
$\Th{Q} \Proves \bforall{x<t}{!A(x)}$.
%
\item O caso do quantificador existencial limitado, no qual temos
!!a{sentence} da forma $\bexists{x < t}{!A(x)}$, é semelhante ao caso
do quantificador universal limitado.
%
\item Suponha que $\lnot \bforall{x<t}{!A(x)}$ seja verdadeira
em~$\Struct{N}$, em que $t$ é um termo fechado. Essa !!{sentence} é
equivalente à !!{sentence} $\bexists{x<t}{\lnot !A(x)}$, sendo a
equivalência derivável em~$\Th{Q}$; portanto, podemos aplicar o
raciocínio para quantificadores existenciais limitados.
%
\item De modo semelhante, suponha que $\lnot \bexists{x<t}!A(x)$ seja
verdadeira em $\Struct{N}$, em que $t$ é um termo fechado. Essa
!!{sentence} é equivalente em $\Th{Q}$ a
$\bforall{x<t}{\lnot!A(x)}$; assim, podemos aplicar o raciocínio para
quantificadores universais limitados.
%
\item Por fim, suponha que $\lnot !A$ seja verdadeira em $\Struct{N}$.
Os únicos casos restantes são aqueles em que $!A$ é atômica e aqueles
em que $\lnot !A \ident \lnot\lnot !B$ para alguma !!{sentence}
$\Delta_0$ $!B$. Se $!A$ é atômica, então, por
\olref{lem:atomic-completeness}, $\Th{Q} \Proves \lnot !A$. Se
$\lnot !A \ident \lnot\lnot !B$, então, pela lógica, ela é
demonstravelmente equivalente em~$\Th{Q}$ a~$!B$, que é verdadeira
em~$\Struct{N}$, pois $\lnot !A$ é verdadeira em~$\Struct{N}$. Pela
hipótese de indução, temos, portanto, $\Th{Q} \Proves \lnot !A$.
\end{enumerate}
\end{proof}

\begin{prob}
Apresente uma prova detalhada do caso existencial em
\olref{lem:delta0-completeness}.
\end{prob}

\begin{thm}
\ollabel{thm:sigma1-completeness}
Se $!A$ é !!a{sentence} $\Sigma_1$ verdadeira em~$\Struct{N}$, então
$\Th{Q} \Proves !A$.
\end{thm}

\begin{proof}
Se $\lexists{x}!A(x)$ é !!a{sentence} $\Sigma_1$ verdadeira
em~$\Struct{N}$, então existem um número natural~$n$ e uma atribuição
de variáveis~$s$ tais que $s(x) = n$ e $\Sat{N}{!A(x)}[s]$. Pelos
fatos usuais sobre a relação de satisfação, segue
$\Sat{N}{!A(\num n)}$. Mas $!A(\num n)$ é uma !!{formula} $\Delta_0$;
portanto, por \olref{lem:delta0-completeness}, temos
$\Th{Q} \Proves !A(\num n)$ e, consequentemente, pela lógica, também
temos $\Th{Q} \Proves \lexists[x][!A(x)]$.
\end{proof}

\end{document}
